\setcounter{equation}{0}
\section{Shor's Factoring Algorithm}
\label{sec:shor}

Shor's algorithm \cite{ptafpf94} is the most famous quantum algorithm: it factors integers in polynomial time, breaking RSA in principle. The full algorithm reduces integer factoring to a quantum subroutine for finding the multiplicative order of a random element, with the order-finding subroutine itself reducing to quantum phase estimation built on the QFT of the previous section.

\subsection{Reduction to order finding}

To factor a composite integer $N$, pick a random $a$ with $\gcd(a,N)=1$. Find the multiplicative order of $a$ modulo $N$ --- the smallest positive integer $r$ with $a^r \equiv 1 \pmod{N}$. Then:
\begin{itemize}
\item if $r$ is odd, restart with a new $a$;
\item if $a^{r/2} \equiv -1 \pmod{N}$, restart;
\item otherwise, $\gcd(a^{r/2} \pm 1,\, N)$ are non-trivial factors of $N$.
\end{itemize}
The probability of obtaining a useful $r$ is bounded below by a constant: Shor showed it is at least $1/2$ for $N$ a product of two distinct odd primes, so a constant number of random choices of $a$ suffice in expectation. The classical post-processing (gcds and the continued-fraction step below) is trivial. The hard part --- the part for which we need a quantum computer --- is finding $r$.

\subsection{Order finding via quantum phase estimation}

Define the unitary $U_a$ on $n = \lceil \log_2 N \rceil$ qubits by
\begin{equation}
U_a \ket{y} \;=\; \ket{a\,y \bmod N}.
\end{equation}
The eigenvalues of $U_a$ are $e^{2\pi i\, s / r}$ for $s = 0,1,\ldots,r-1$. Quantum phase estimation extracts a (good rational approximation to a) random one of these eigenvalues, and the continued-fraction expansion recovers the denominator $r$.

The circuit uses two registers:
\begin{itemize}
\item a \emph{counting} register of $t = 2n+1$ qubits, initialised to $\ket{0}^{\otimes t}$;
\item a \emph{target} register of $n$ qubits, initialised to $\ket{1}$.
\end{itemize}
The algorithm proceeds:
\begin{enumerate}
\item Apply $H^{\otimes t}$ to the counting register, producing the uniform superposition.
\item For each counting qubit $j = 0,1,\ldots,t-1$, apply a controlled-$U_a^{2^j}$ from counting qubit $j$ to the target register. This entangles the counting register with $\ket{a^x \bmod N}$ in the target.
\item Apply the inverse QFT (Section~\ref{sec:qft}) to the counting register.
\item Measure the counting register to obtain an integer $c \in \{0,1,\ldots,2^t-1\}$.
\item Compute $r$ from $c/2^t$ via the continued-fraction expansion.
\end{enumerate}
After step 2 the joint state is
\begin{equation}
\frac{1}{\sqrt{2^t}}\sum_{x=0}^{2^t - 1} \ket{x} \ket{a^x \bmod N}.
\end{equation}
The inverse QFT in step 3 concentrates probability mass on values of $c$ such that $c/2^t \approx s/r$ for some integer $s\in\{0,1,\ldots,r-1\}$, and the continued-fraction algorithm recovers $r$ from this rational approximation with high probability.

\subsection{The expensive subroutine: modular exponentiation}

Implementing the controlled-$U_a^{2^j}$ gates is the dominant cost of Shor's algorithm. The standard construction builds quantum ripple-carry adders, lifts those to modular addition, lifts modular addition to modular multiplication, and finally lifts modular multiplication to modular exponentiation. The most efficient known circuits use $O(n^3)$ gates and dominate the algorithm's total gate count. The Beauregard $2n+1$-qubit circuit \cite{beauregard2003} is the standard textbook construction.

In a state-vector simulator, there are two options.

\paragraph{Honest implementation.}
Build the modular exponentiation circuit from the same quantum adder primitives a physical machine would use. Useful for studying noise, gate counts, and ancilla overhead. The Beauregard circuit and its variants take this route.

\paragraph{Simulator shortcut.} Since the simulator has direct access to all amplitudes, we can update the target register's basis-state index classically for each non-zero amplitude in the counting register. The simulated quantum state is identical to the one produced by the honest circuit, but the time and ancilla costs are dramatically lower:

\begin{lstlisting}[caption={Modular exponentiation, single-process simulator shortcut. Iterates over basis states, classically computes $a^x \bmod N$, and rewrites the target register. See note below for the distributed case.}]
/* Note: this version assumes a single-process simulator: new_global may
 * fall on a different rank than idx, but local_of() here treats the
 * whole state vector as local. A distributed version must bucket the
 * (new_global, amp[idx]) pairs by destination rank, exchange them with
 * an MPI_Alltoallv, then accumulate into new_amp on each rank.       */
/* mod_pow: uint64_t mod_pow(uint64_t base, uint64_t exp, uint64_t mod); */
void apply_modular_exp(qreg *q, int counting_start, int t,
                       int target_start, int n,
                       uint64_t a, uint64_t N) {
    complex double *new_amp = calloc(q->local_size, sizeof *new_amp);
    for (size_t idx = 0; idx < q->local_size; idx++) {
        if (cabs(q->amp[idx]) < 1e-15) continue;
        size_t   global = global_index(q, idx);
        uint64_t x      = (global >> counting_start) & ((1ULL << t) - 1);
        uint64_t y      = (global >> target_start)   & ((1ULL << n) - 1);
        uint64_t y_new  = (y * mod_pow(a, x, N)) % N;
        size_t   new_global = (global
            & ~(((1ULL << n) - 1) << target_start))
            | ((size_t)y_new << target_start);
        new_amp[local_of(q, new_global)] += q->amp[idx];
    }
    memcpy(q->amp, new_amp, q->local_size * sizeof *new_amp);
    free(new_amp);
}
\end{lstlisting}

The listing as written assumes the entire state vector is held on a single process; under the distributed layout of Section~\ref{sec:parallel}, $U_a^x$ permutes basis-state indices in a way that can move amplitude across rank boundaries, so a distributed implementation must perform an all-to-all redistribution of $(\text{new\_global}, q\!\to\!\text{amp}[\text{idx}])$ pairs --- typically by bucketing per destination rank and exchanging with \verb|MPI_Alltoallv| --- before accumulating into \verb|new_amp|. This shortcut is dishonest as a model of physical execution --- it tells us nothing about real gate count or noise --- but it produces the correct quantum state for studying the information-theoretic behaviour of the algorithm, and runs in seconds where the honest circuit would take hours.

\subsection{Classical post-processing}

After measurement we have $c \in \{0, \ldots, 2^t - 1\}$. Compute the continued-fraction expansion of $c/2^t$ and find the convergent $p/q$ with $q < N$ and $|c/2^t - p/q| < 1/(2 \cdot 2^t)$. If $q$ is the true order of $a$ modulo $N$, we are done; otherwise verify $a^q \equiv 1 \pmod{N}$, and if not, restart the quantum portion of the algorithm. The continued-fraction step succeeds with constant probability per quantum run, so a constant expected number of repetitions suffices.

\subsection{Resource requirements}

To factor an $n$-bit integer $N$, Shor's algorithm requires approximately $2n+1$ counting qubits plus $n$ target qubits plus $\Theta(n)$ ancillas for modular arithmetic --- roughly $5n$ qubits in total. For RSA-2048 this is about $10{,}000$ logical qubits; under realistic surface-code overhead the physical qubit count rises to millions. This explains why factoring real-world RSA keys remains out of reach despite the algorithm being more than three decades old.

In simulation, the qubit count is the binding constraint, not the gate count, and the exact count depends on whether one simulates the \emph{honest} modular-arithmetic circuit (counting + target + arithmetic ancillas, $\approx 5n$ qubits in total) or applies a \emph{shortcut} that replaces the modular-exponentiation subcircuit with an in-place permutation of state-vector amplitudes, giving $\approx 3n+1$ qubits (counting register $t = 2n+1$ plus target register $n$, no arithmetic ancillas required). The bundled \texttt{implementation/c/}, \texttt{implementation/go/}, and \texttt{implementation/python/} libraries all use the shortcut: \texttt{apply\_modular\_exp} computes $\ket{x}\ket{y} \mapsto \ket{x}\ket{y\cdot a^x \bmod N}$ classically and then writes the result back as a single bulk permutation of the amplitude array. This is faithful to the abstract specification of the gate but is not a quantum circuit; an honest construction would unroll it into $O(n^3)$ elementary gates with $\Theta(n)$ ancillas \cite{beauregard2003}.
\begin{itemize}
\item Honest-circuit model: factoring $N=15$ ($n=4$) requires roughly $5n=20$ qubits ($16\,\text{MB}$ of state-vector memory --- fits on a phone). Factoring $N=21$ ($n=5$) requires $\sim 25$ qubits ($512\,\text{MB}$ --- fits on a laptop). Factoring a 32-bit number requires $\sim 160$ qubits, far beyond what state-vector simulation on \emph{any} conceivable classical hardware can reach.
\item Shortcut model (as bundled): factoring $N=15$ requires $3n+1 = 13$ qubits ($n=4$ target bits plus $t = 2n+1 = 9$ counting bits). Factoring $N=21$ requires $3n+1 = 16$ qubits; this is the configuration exercised by the \texttt{RUN\_SHOR\_21=1} gated test in all three implementations.
\end{itemize}
Either way, the simulation budget grows exponentially with $n$ in memory, which is the asymptotic gap that motivated the entire field --- and the reason why simulation is fundamentally a tool for algorithm design and small-scale validation, not a substitute for a physical quantum computer.
